\begin{titlepage}
\centering

\includegraphics[height=1.4in]{figuras/escudo-unap.png}\par\vspace{0.8em}

{\Large\textbf{UNIVERSIDAD NACIONAL DEL ALTIPLANO - PUNO}}\par\vspace{0.4em}
{\large\textbf{FACULTAD DE INGENIERÍA MECÁNICA ELÉCTRICA, ELECTRÓNICA Y SISTEMAS}}\par\vspace{0.3em}
{\normalsize\textbf{ESCUELA PROFESIONAL DE INGENIERÍA MECÁNICA ELÉCTRICA}}\par\vspace{1.8em}

\textcolor{modelblue}{\rule{0.85\textwidth}{1.8pt}}\par\vspace{1.2em}

{\LARGE\bfseries\color{headingred} EXPEDIENTE TÉCNICO DE INGENIERÍA}\par\vspace{1.0em}

{\Large\bfseries PROYECTO:}\par\vspace{0.5em}
{\Large\textbf{``ELECTRIFICACIÓN Y REDES DE DISTRIBUCIÓN EN MEDIA TENSIÓN (10 kV) Y BAJA TENSIÓN (380/220 V) DESDE LA MINICENTRAL HIDROELÉCTRICA LIMBANI (200 kW, TURBINA MICHELL-BANKI)''}}\par\vspace{1.2em}

{\large\textbf{CASCO URBANO DEL DISTRITO DE LIMBANI -- PROVINCIA DE SANDIA -- DEPARTAMENTO DE PUNO}}\par\vspace{1.2em}

\textcolor{modelblue}{\rule{0.85\textwidth}{1.8pt}}\par\vspace{2.2em}

\begin{flushleft}
\large
\textbf{CURSO:} Instalaciones Eléctricas II / Redes Eléctricas y Subestaciones\par\vspace{0.4em}
\textbf{DOCENTE:} Mg. Gregorio Meza Marocho\par\vspace{0.4em}
\textbf{RESPONSABLE / AUTOR:} Br. Aquiles Taylor Ramos Yapo\par\vspace{0.4em}
\textbf{SISTEMA GEODÉSICO:} WGS 84 / UTM Zona 19 Sur (EPSG:32719)\par\vspace{0.4em}
\textbf{POTENCIA DE GENERACIÓN:} 200 kW (Turbina Michell-Banki / Alternador 380 V)\par\vspace{0.4em}
\textbf{PREDIOS / VIVIENDAS BENEFICIADAS:} 764 Viviendas Consolidadas en Catastro Oficial
\end{flushleft}

\vfill

{\large\bfseries PUNO -- PERÚ}\par
{\normalsize 2026}

\end{titlepage}
